\chapter{Contexte du projet et fondements techniques}
\label{chap:contexte}

\section*{Introduction}

Ce premier chapitre situe le stage dans son cadre : l'organisme d'accueil, les objectifs et le
périmètre volontairement borné du projet. Il présente ensuite les fondements techniques sur
lesquels repose le système --- et en particulier la raison pour laquelle un grand modèle de
langage seul ne peut pas répondre de façon fiable à une question juridique.

\section{Organisme d'accueil}

Pulsaride Solutions est une entreprise du domaine \textit{deeptech}, spécialisée en sécurité et
en intelligence artificielle, dont le siège est situé rue Henri Dunant à Pessac, en France.

Le stage s'est inscrit dans l'axe « intelligence artificielle appliquée » de l'entreprise, sur un
sujet articulant traitement automatique du langage et souveraineté des données. L'encadrement a
été assuré par M. Lyazid \textsc{Salihi}.

\section{Objectifs et missions du stage}

L'objectif est de concevoir et d'évaluer un assistant de recherche juridique souverain fondé sur
l'architecture RAG. Le citoyen pose une question en langage naturel ; le système recherche le
passage pertinent dans le corpus juridique, puis formule une réponse claire \textbf{en citant
systématiquement la source officielle}.

Le principe de fiabilité est central : le système s'appuie sur les textes réels et
\textbf{indique honnêtement lorsqu'il ne dispose pas de l'information}, plutôt que d'inventer.

\begin{tcolorbox}[colback=gray!5,colframe=gray!40,title=Exemple d'interaction visée]
\textbf{Citoyen :} « J'ai droit à combien de jours de congé de maternité ? »

\medskip
\textbf{Assistant :} « Selon l'article 152 du Code du travail (Loi 65-99), la salariée bénéficie
d'un congé de maternité de 14 semaines. Le contrat de travail est suspendu, non rompu, pendant
cette période.

\textit{Source : Code du travail, article 152.}

Réponse fournie à titre informatif ; pour une situation particulière, consultez un
professionnel. »
\end{tcolorbox}

Le stage s'articule autour de quatre missions :

\begin{enumerate}
    \item \textbf{Constitution et préparation du corpus juridique.} Collecter, nettoyer et
    découper intelligemment le corpus par article. La qualité de ce découpage conditionne
    directement la qualité des réponses.
    \item \textbf{Moteur de recherche sémantique et génération.} Indexer le corpus, retrouver les
    passages pertinents à partir d'une question, puis générer une réponse claire et sourcée, en
    garantissant l'ancrage dans les textes réels et la gestion explicite du « je ne sais pas ».
    \item \textbf{Évaluation de la fiabilité.} Construire un jeu de questions-réponses de
    référence et mesurer la précision du système. L'analyse honnête des limites fait partie
    intégrante du livrable.
    \item \textbf{Extension multilingue} \textit{(ambition complémentaire)}. Explorer une couche
    de normalisation linguistique permettant de comprendre une question posée en arabe ou en
    darija, tout en conservant l'ancrage des réponses dans le texte officiel d'origine.
\end{enumerate}

\section{Périmètre et non-objectifs}

Pour garantir la faisabilité, le sujet a été volontairement borné : un \textbf{périmètre
juridique unique} --- le droit du travail --- et un fonctionnement \textbf{monolingue solide} en
français. L'extension multilingue et l'élargissement à d'autres codes constituent des ambitions
complémentaires, non des engagements.

Le cadrage négatif est aussi important que le cadrage positif. Le système n'est explicitement
pas :

\begin{itemize}
    \item \textbf{un outil de conseil juridique personnalisé.} Il relève de la \emph{restitution
    d'information juridique} ; chaque réponse porte un avertissement en ce sens ;
    \item \textbf{un système qui répond « au mieux ».} Face à une question sans réponse dans le
    corpus, \textbf{l'abstention est le comportement correct}, pas un échec ;
    \item \textbf{un produit fini.} Le livrable est le moteur RAG et son évaluation.
\end{itemize}

Quatre interdictions ont été posées dès la conception et tenues jusqu'au terme du stage : pas de
\textit{fine-tuning} du modèle, pas de mémoire conversationnelle multi-tours, pas d'autre code
juridique que le Code du travail, et aucune intégration d'API externe dans le cœur du système.

\section{Organisation du travail}

Le travail a été organisé en huit phases, chacune assortie d'un livrable et d'un risque principal
identifié à l'avance. Un \textbf{journal de bord} a été tenu quotidiennement, consignant chaque
décision, chaque résultat et chaque difficulté rencontrée --- y compris les impasses. Ce journal
a servi de matière première à la rédaction du présent rapport.

\begin{table}[H]
    \centering
    \begin{tabular}{|m{1.2cm}|m{6.5cm}|m{6.5cm}|}
        \hline
        \textbf{Phase} & \textbf{Objet} & \textbf{Livrable} \\
        \hline
        S1 & Acquisition du corpus, extraction, découpage par article & Liste d'articles brute, fidélité vérifiée \\
        \hline
        S2 & Structuration, métadonnées, hiérarchie & Corpus reproductible au format JSONL \\
        \hline
        S3 & Premier pipeline RAG bout-en-bout & CLI : question $\rightarrow$ réponse citée \\
        \hline
        S4 & Ancrage, citation, abstention & Réponses fiables, abstention hors périmètre \\
        \hline
        S5 & Recherche hybride, jeu de référence & Recherche améliorée, jeu d'évaluation \\
        \hline
        S6 & Exécution de l'évaluation et métriques & Rapport de performance, analyse d'erreurs \\
        \hline
        S7 & Durcissement, reproductibilité & Prototype propre et reproductible \\
        \hline
        S8 & Rapport et soutenance & Rapport, démonstration, dépôt reproductible \\
        \hline
    \end{tabular}
    \caption{Découpage du stage en phases}
    \label{tab:phases}
\end{table}

\section{Fondements techniques}

\subsection{Grands modèles de langage et hallucination}

Un grand modèle de langage (\textit{Large Language Model}, LLM) est un modèle statistique
entraîné à prédire la suite la plus plausible d'une séquence de texte. Cette formulation explique
à la fois sa fluidité et sa principale faiblesse : le modèle produit ce qui est
\emph{vraisemblable}, non ce qui est \emph{vrai}.

Ce comportement porte le nom d'\textbf{hallucination}. Dans un contexte juridique, il est
particulièrement dangereux : un modèle sollicité sur un article de loi qu'il ne connaît pas
produira volontiers un numéro d'article plausible accompagné d'un contenu inventé. La réponse
paraît crédible, la référence semble vérifiable, et rien dans la forme ne signale l'invention.

Deux stratégies existent pour ancrer un modèle dans des faits réels : le
\textbf{\textit{fine-tuning}}, coûteux, difficile à mettre à jour, et qui ne permet toujours pas
de citer une source vérifiable ; et la \textbf{génération augmentée par la recherche} (RAG), qui
fournit au modèle les documents pertinents \emph{au moment de la question} et lui demande de
s'appuyer uniquement sur eux. C'est la seconde qui a été retenue : elle permet la citation, donc
la vérification, condition même de la fiabilité recherchée.

\subsection{L'architecture RAG}

Le principe du RAG consiste à séparer deux préoccupations habituellement confondues :
\emph{savoir} et \emph{formuler}. Le savoir est stocké dans un corpus documentaire externe et
interrogeable ; le modèle de langage n'est plus qu'un moteur de reformulation appliqué à des
documents qu'on lui fournit.

Le système se décompose en deux pipelines : un pipeline d'\textbf{indexation}, exécuté une seule
fois hors ligne, et un pipeline de \textbf{requête}, exécuté à chaque question. Cette séparation
a une conséquence pratique décisive pour le débogage : il est possible d'inspecter les documents
retrouvés \emph{avant} l'appel au modèle. Or la majorité des mauvaises réponses en RAG
proviennent d'une mauvaise \emph{recherche}, non d'une mauvaise \emph{génération}.

\subsection{Recherche sémantique par embeddings}

Un \textbf{embedding} est la représentation d'un texte sous forme de vecteur numérique de
dimension fixe, construite de telle sorte que deux textes de sens proche produisent des vecteurs
proches dans l'espace. La proximité est mesurée par la \textbf{similarité cosinus} :

\begin{equation}
    \mathrm{sim}(\vec{a}, \vec{b}) = \frac{\vec{a} \cdot \vec{b}}{\|\vec{a}\|\,\|\vec{b}\|}
\end{equation}

Pour des vecteurs normalisés, cette expression se réduit à un produit scalaire, borné entre $-1$
et $1$.

L'intérêt de cette approche est qu'elle capture le \emph{sens} et non les mots : une question
posée en langage courant (« mon patron peut me virer sans raison ? ») peut retrouver un article
rédigé en langage juridique formel, sans partager aucun terme avec lui. Sa limite est
symétrique : elle est moins précise lorsque le terme exact compte --- un numéro d'article, une
expression figée comme « Loi 65-99 ».

\subsection{Recherche lexicale : BM25}

BM25 (\textit{Best Matching 25}) est une fonction de classement lexicale qui score un document en
fonction des termes qu'il partage avec la requête. Elle pondère chaque terme par sa rareté dans
le corpus --- un mot rare est plus discriminant qu'un mot fréquent --- et normalise par la
longueur du document. BM25 est exactement complémentaire des embeddings : excellent sur les
correspondances exactes, aveugle à la reformulation.

\subsection{Fusion des classements : Reciprocal Rank Fusion}

Combiner les deux méthodes pose un problème pratique : leurs scores ne sont pas comparables. Un
score BM25 n'est pas borné ; une similarité cosinus est bornée entre $-1$ et $1$. Faire la
moyenne de $8$ et de $0{,}7$ n'a aucun sens.

La \textbf{fusion par rang réciproque} (\textit{Reciprocal Rank Fusion}, RRF) contourne ce
problème en n'utilisant que le \emph{rang} de chaque document dans chaque liste :

\begin{equation}
    \mathrm{RRF}(d) = \sum_{m \in M} \frac{1}{k + \mathrm{rang}_m(d)}
\end{equation}

où $M$ est l'ensemble des méthodes de recherche, $\mathrm{rang}_m(d)$ le rang du document $d$
selon la méthode $m$, et $k$ une constante d'amortissement (valeur usuelle : 60). Un document
bien classé par les deux méthodes obtient un score élevé, sans qu'aucune normalisation d'échelle
ne soit nécessaire.

\subsection{La contrainte d'exécution locale}

La contrainte de souveraineté a une conséquence technique directe : \textbf{aucun composant du
cœur du système ne peut être un service en ligne}. Cela exclut les embeddings par API, les
services de reclassement commerciaux et les modèles de génération propriétaires accessibles
uniquement à distance. Toute la chaîne repose donc sur des modèles open-source exécutables sur du
matériel ordinaire. Cette contrainte, loin d'être une simple limitation, structure l'ensemble des
choix présentés au chapitre suivant.

\newpage
\section*{Conclusion}

Le projet répond à un besoin concret --- rendre le droit du travail interrogeable en langage
naturel --- sous une contrainte structurante, la souveraineté des données, et avec une exigence
centrale, la fiabilité. Le RAG apporte une réponse directe au problème de l'invention de
références : en fournissant au modèle les textes au moment de la question, il rend la citation
possible et la vérification réalisable. Le chapitre suivant traduit ces principes en une
architecture concrète.
